\SourceLeaf{283}%
\[
\begin{aligned}
l&=l^0+\lambda x+\lambda'xx+\text{usw.}\\
m&=m^0(1+\mu x+\mu'xx+\text{usw.})
\end{aligned}
\]
entwickelt sind, so ergibt die Rechnung
\[
\begin{aligned}
\operatorname{tang}\psi={}&(1+\mu x+\mu'xx+\text{usw.})\operatorname{tang}\psi^0\\
&-l^0x-\tfrac{1}{2}(\lambda+l^0\mu)xx\\
&-(\tfrac{1}{3}\lambda'+\tfrac{1}{6}\lambda\mu-\tfrac{1}{6}l^0\mu\mu+\tfrac{2}{3}l^0\mu')x^3-\text{usw.}
\end{aligned}
\]
und daraus, weil \(u=\int\operatorname{tang}\psi\mathord{.}\mathrm{d}x\),
\[
\begin{aligned}
u={}&(x+\tfrac{1}{2}\mu xx+\tfrac{1}{3}\mu'x^3+\text{usw.})\operatorname{tang}\psi^0\\
&-\tfrac{1}{2}l^0xx-\tfrac{1}{6}(\lambda+l^0\mu)x^3\\
&-(\tfrac{1}{12}\lambda'+\tfrac{1}{24}\lambda\mu-\tfrac{1}{24}l^0\mu\mu+\tfrac{1}{6}l^0\mu')x^4-\text{usw.}
\end{aligned}
\]
wobei keine Konstante hinzuzufügen ist, weil für \(x=0\) auch \(u=0\) gilt. Da nun auch für \(x=h\) die Gleichung \(u=0\) gilt, folgt aus dieser Gleichung
\[
\operatorname{tang}\psi^0=\tfrac{1}{2}l^0h+(\tfrac{1}{6}\lambda-\tfrac{1}{12}l^0\mu)hh+(\tfrac{1}{12}\lambda'-\tfrac{1}{24}\lambda\mu)h^3+\text{usw.}
\]
Setzt man in der Gleichung für \(\psi\) auch für \(x\) den Wert \(h\) und für \(\operatorname{tang}\psi^0\) den eben gefundenen Wert ein, so ergibt sich
\[
\begin{aligned}
\operatorname{tang}\psi'={}&-\tfrac{1}{2}l^0h-(\tfrac{1}{3}\lambda+\tfrac{1}{12}l^0\mu)hh\\
&-(\tfrac{1}{4}\lambda'+\tfrac{1}{24}\lambda\mu-\tfrac{1}{12}l^0\mu\mu+\tfrac{1}{6}l^0\mu')h^3\quad\text{usw.}
\end{aligned}
\]
\Needspace{19\baselineskip}
Da
\[
\begin{aligned}
l'&=l^0+\lambda h+\lambda'hh+\text{usw}\\
m'&=m^0(1+\mu h+\mu'hh+\text{usw.})
\end{aligned}
\]
gilt, erhält man
\[
\begin{aligned}
(\tfrac{1}{3}l^0+\tfrac{1}{6}l')h\BareRoot[6]{\frac{m^0}{m'}}={}&\tfrac{1}{2}l^0h+(\tfrac{1}{6}\lambda-\tfrac{1}{12}l^0\mu)hh\\
&+(\tfrac{1}{6}\lambda'-\tfrac{1}{36}\lambda\mu+\tfrac{7}{144}l^0\mu\mu-\tfrac{1}{12}l^0\mu')h^3\quad\text{usw.}
\end{aligned}
\]
\[
\begin{aligned}
-(\tfrac{1}{6}l^0+\tfrac{1}{3}l')h\BareRoot[6]{\frac{m'}{m^0}}={}&-\tfrac{1}{2}l^0h-(\tfrac{1}{3}\lambda+\tfrac{1}{12}l^0\mu)hh\\
&-(\tfrac{1}{3}\lambda'+\tfrac{1}{18}\lambda\mu-\tfrac{5}{144}l^0\mu\mu+\tfrac{1}{12}l^0\mu')h^3\quad\text{usw.}
\end{aligned}
\]
Diese Ausdrücke stimmen also in den ersten beiden Gliedern, das heißt bis zur Ordnung \(hh\), mit den obigen Werten von \(\operatorname{tang}\psi^0,\operatorname{tang}\psi'\) überein. Man kann diese bequemen Ausdrücke daher als hinreichend genaue Werte der Tangenten oder unter Hinzunahme des Faktors \(206265''\) als Werte der Winkel \(\psi^0,\psi'\) selbst annehmen.

Die Länge der Linie \(L\) selbst zwischen den Punkten auf dem Ellipsoid, denen auf der Kugel die Punkte \(F,G\) entsprechen, ist das Integral
